\section{量子ビットと量子状態}
\label{sec:qubit-state}

\subsection{ブラケット記法}

複素数を成分にもつ有限次元のベクトル空間を$\mathcal H$とする．
空間$\mathcal H$には内積が定められているものとする．
$\mathcal H$のベクトルを
\begin{equation}
  \ket{\psi}
  \coloneq
  \begin{pmatrix}
    \psi_0&\psi_1&\cdots&\psi_{d-1}
  \end{pmatrix}^{\mathsf T}
  \label{eq:ket-coordinate}
\end{equation}
と書き，これを\textbf{ケット}という．
ケット$\ket{\psi}$の各成分を複素共役にして横に並べたものを
$\bra{\psi}:=\ket{\psi}^{\dagger}$と書き，\textbf{ブラ}という．
ベクトル$\ket{\psi}$と$\ket{\phi}$の内積は
$\braket{\phi|\psi}:=\bra{\phi}\ket{\psi}$と書く．座標を使えば，
\begin{equation}
  \braket{\phi|\psi}=\sum_{j=0}^{d-1}\overline{\phi_j}\psi_j
\end{equation}
である．内積から定まるノルムを
$\lVert\ket{\psi}\rVert:=\sqrt{\braket{\psi|\psi}}$と書く．
また，$\ket{\psi}\bra{\phi}$は線形写像を表す．
この写像は$\ket{v}$を$\braket{\phi|v}\ket{\psi}$へ移す．

% \begin{prop}
% 正規直交基底$\{\ket{e_j}\}_{j=0}^{d-1}$に対し，任意の$\ket{\psi}\in\mathcal H$は
% \begin{equation}
%   \ket{\psi}=\sum_{j=0}^{d-1}\braket{e_j|\psi}\ket{e_j}
%   \label{eq:orthonormal-expansion}
% \end{equation}
% と一意的に表される．
% \end{prop}
% \begin{proof}
% 基底であるから$\ket{\psi}=\sum_j c_j\ket{e_j}$と一意的に表される．
% 左から$\bra{e_k}$を作用させると
% \[
%   \braket{e_k|\psi}
%   =\sum_jc_j\braket{e_k|e_j}
%   =\sum_jc_j\delta_{kj}=c_k
% \]
% である．したがって$c_j=\braket{e_j|\psi}$となる．
% \end{proof}

\subsection{量子状態の公理}

% \begin{defi}[純粋状態]
% ほかの系から影響を受けない量子系を考える．
% その純粋状態は，状態空間$\mathcal H$の長さ1のベクトル$\ket{\psi}$で表す．
% つまり，$\braket{\psi|\psi}=1$とする．
% \end{defi}

ベクトルをその長さで割り，長さを1にそろえることを\textbf{規格化}という．

1量子ビットの状態空間は$\mathbb C^2$である．計算基底を
\begin{equation}
  \ket{0}=\begin{pmatrix}1\\0\end{pmatrix},
  \qquad
  \ket{1}=\begin{pmatrix}0\\1\end{pmatrix}
\end{equation}
と定めると，一般の1量子ビット状態は
\begin{equation}
  \ket{\psi}=\alpha\ket{0}+\beta\ket{1},
  \qquad |\alpha|^2+|\beta|^2=1
  \label{eq:single-qubit-state}
\end{equation}
と表される．係数$\alpha,\beta$を\textbf{確率振幅}という．
実数$\gamma$に対して$\ket{\psi}$と$e^{\mathrm{i}\gamma}\ket{\psi}$は
同じ物理状態を表す．この因子を\textbf{グローバル位相}という．
これに対して，$\alpha$と$\beta$の間の位相の差は，後の演算と測定結果に影響する．
この差を\textbf{相対位相}という．

\subsection{射影測定の公理}

\begin{defi}[直交射影]
線形写像$P:\mathcal H\to\mathcal H$が$P^2=P$かつ$P^\dagger=P$を満たすとき，
$P$を直交射影という．
\end{defi}

\begin{defi}[射影測定]
測定結果の集合を$M$とする．各結果$m\in M$に対して直交射影$P_m$を用意し，
\begin{equation}
  P_mP_{m'}=\delta_{mm'}P_m,
  \qquad \sum_{m\in M}P_m=I
  \label{eq:projective-measurement-completeness}
\end{equation}
を満たすとする．状態$\ket{\psi}$の測定で結果$m$を得る確率は
\begin{equation}
  p(m)=\bra{\psi}P_m\ket{\psi}
      =\lVert P_m\ket{\psi}\rVert^2
  \label{eq:born-rule-projector}
\end{equation}
である．この確率をボルン則という．確率が$p(m)>0$なら，測定後の状態は
\begin{equation}
  \ket{\psi_m}=\frac{P_m\ket{\psi}}{\sqrt{p(m)}}
  \label{eq:post-measurement-state}
\end{equation}
であるとする．
\end{defi}

\begin{prop}
射影測定の各結果の確率は非負で総和は1であり，
式\eqref{eq:post-measurement-state}の状態は規格化されている．
\end{prop}
\begin{proof}
式\eqref{eq:born-rule-projector}はノルムの二乗であるから$p(m)\geq0$である．
完全性関係と状態の規格化より
\[
  \sum_m p(m)
  =\bra{\psi}\left(\sum_mP_m\right)\ket{\psi}
  =\bra{\psi}I\ket{\psi}=1
\]
を得る．さらに$P_m^\dagger P_m=P_m$より
\[
  \braket{\psi_m|\psi_m}
  =\frac{\bra{\psi}P_m^\dagger P_m\ket{\psi}}{p(m)}=1
\]
となる．
\end{proof}

計算基底で測定する場合は，$P_0=\ket{0}\bra{0}$，$P_1=\ket{1}\bra{1}$とする．
式\eqref{eq:single-qubit-state}に対し，$p(0)=|\alpha|^2$，
$p(1)=|\beta|^2$である．

\begin{prop}
状態$\ket{\psi}$と$e^{\mathrm{i}\gamma}\ket{\psi}$は，すべての射影測定で同じ確率を与える．
\end{prop}
\begin{proof}
任意の$m$について
\[
  e^{-\mathrm{i}\gamma}\bra{\psi}P_m
  e^{\mathrm{i}\gamma}\ket{\psi}
  =\bra{\psi}P_m\ket{\psi}
\]
であるため，測定確率は大域位相に依存しない．
\end{proof}

\begin{ex}
状態$\ket{\psi}=(\ket{0}+\mathrm{i}\ket{1})/\sqrt2$を計算基底で測定すると，
結果$0$と$1$はともに確率$1/2$で得られる．
\end{ex}
